\section{幻灯片与可视化证据}

\subsection{策略与证据分层：Slide 0-04}

\begin{figure}[H]
\centering
\includegraphics[width=0.8\textwidth]{slide-0-04.jpg}
\caption{Slide 0-04 把 reward model、policy update、治理流程叠加在一张画布上，强调 evidence-to-action 的 track。}
\end{figure}

Slide 0-04 展示了多个 evidence signal 如何在同一个 timeline 中被收集与审查，提醒 teams 不能只关注一个 benchmark，而要将 evidence matrix 变成操作性的 dashboard。

\begin{knowledgebox}{分层 evidence dashboard}
在 dashboard 中将 benchmark、human eval、deployment logs 做 stacked view，可以更容易定位哪一层 signal 首先 drift。
\end{knowledgebox}

\subsection{运营节奏：Slide 0-05}

\begin{figure}[H]
\centering
\includegraphics[width=0.8\textwidth]{slide-0-05.jpg}
\caption{Slide 0-05 画出了 alignment-engineering 的节奏感（data → training → deployment → monitoring），结合 action items。}
\end{figure}

每个阶段都配有 artifacts：data stage 有 prompt catalog；training stage 有 KL margin report；deployment stage 有 drift alarm；monitoring stage 有 ops playbook。这样的节奏图可以作为 quarterly planning 的 reference。

\begin{importantbox}{操作节奏的关键}
用 timeline+artifact view 让跨职能团队可以在 planning 会议中快速达成共识，从而避免 alignment 逻辑在交付与监控之间断层。
\end{importantbox}

\subsection{质量矩阵：Slide 0-06}

\begin{figure}[H]
\centering
\includegraphics[width=0.8\textwidth]{slide-0-06.jpg}
\caption{Slide 0-06 把质量、信任与治理矩阵化，强调不同 stakeholder 的衡量标准。}
\end{figure}

该幻灯片把 quality gate、trust signal 和 governance checkpoints 组合成一个矩阵，便于把 evidence 逐层传递给 product、ops、legal。

\begin{warningbox}{不要孤立强化某条 signal}
只加强 automation metric（如 latency）而不结合 governance checklist，会导致高效但不安全的版本上线。Slide 0-06 提醒我们要在矩阵中保持 balance。
\end{warningbox}

\subsection{本章小结}

这些幻灯片提供了 visual anchor，把 abstract 的 RLHF pipeline 编织成 layer-by-layer 的 evidence-dashboard、ops rhythm、governance matrix，有助于 cross-team alignment。

